\documentclass[12pt,a4paper]{article}

% --- Metadata ---
\def\courseshortheader{HỌC MÁY --- GIỚI THIỆU}
\def\coverbookline{TRÍ TUỆ NHÂN TẠO}
\def\covermaintitle{Giới thiệu học máy}
\def\coversubtitle{Lý thuyết --- Khái niệm nền, học có giám sát}

% Shared preamble for nhom_5 (requires \courseshortheader before \input)
\usepackage[utf8]{inputenc}
\usepackage[vietnamese]{babel}
\usepackage{amsmath,amssymb,amsthm}
\usepackage{geometry}
\usepackage{enumitem}
\usepackage[most]{tcolorbox}
\usepackage{hyperref}
\usepackage{fancyhdr}
\usepackage{graphicx}
\usepackage{booktabs}
\usepackage{tikz}
\usetikzlibrary{calc,arrows.meta,decorations.pathreplacing,positioning}
\usepackage{eso-pic}
\usepackage{xcolor}

\geometry{a4paper, margin=2cm, bottom=2.5cm}
\hypersetup{colorlinks=true, linkcolor=blue!70!black, urlcolor=blue}
\setlist[itemize]{topsep=4pt, itemsep=2pt}
\setlist[enumerate]{topsep=4pt, itemsep=2pt}

\AddToShipoutPictureFG{%
  \AtPageCenter{%
    \begin{tikzpicture}[remember picture, overlay]
      \node[rotate=45, scale=1, opacity=0.06]{%
        \fontsize{60}{72}\selectfont\bfseries\sffamily Đặng Tấn Quí%
      };
    \end{tikzpicture}%
  }%
}

\pagestyle{fancy}
\fancyhf{}
\fancyhead[L]{\small\textit{\courseshortheader}}
\fancyhead[R]{\small\textit{Đặng Tấn Quí}}
\fancyfoot[C]{\thepage}
\fancyfoot[L]{\footnotesize\textcopyright\ Đặng Tấn Quí}
\fancyfoot[R]{\footnotesize SĐT: 0377\,122\,210}
\renewcommand{\headrulewidth}{0.4pt}
\renewcommand{\footrulewidth}{0.4pt}

\fancypagestyle{plain}{%
  \fancyhf{}
  \fancyfoot[C]{\thepage}
  \fancyfoot[L]{\footnotesize\textcopyright\ Đặng Tấn Quí}
  \fancyfoot[R]{\footnotesize SĐT: 0377\,122\,210}
  \renewcommand{\headrulewidth}{0pt}
  \renewcommand{\footrulewidth}{0.4pt}
}

\newtcolorbox{dongco}[1][]{%
  enhanced, breakable,
  colback=teal!4!white, colframe=teal!60!black,
  fonttitle=\bfseries, title={Động cơ học -- Tại sao cần khái niệm này?}, #1%
}
\newtcolorbox{ynghia}[1][]{%
  enhanced, breakable,
  colback=purple!3!white, colframe=purple!50!black,
  fonttitle=\bfseries, title={Ý nghĩa \& Trực giác}, #1%
}
\newtcolorbox{ungdung}[1][]{%
  enhanced, breakable,
  colback=olive!5!white, colframe=olive!60!black,
  fonttitle=\bfseries, title={Ứng dụng thực tế}, #1%
}
\newtcolorbox{dinhnghia}[1][]{%
  enhanced, breakable,
  colback=blue!3!white, colframe=blue!60!black,
  fonttitle=\bfseries, title={Định nghĩa}, #1%
}
\newtcolorbox{dinhly}[1][]{%
  enhanced, breakable,
  colback=green!3!white, colframe=green!50!black,
  fonttitle=\bfseries, title={Định lý}, #1%
}
\newtcolorbox{tinhchat}[1][]{%
  enhanced, breakable,
  colback=yellow!5!white, colframe=orange!50!black,
  fonttitle=\bfseries, title={Tính chất}, #1%
}
\newtcolorbox{phuongphap}[1][]{%
  enhanced, breakable,
  colback=gray!5!white, colframe=gray!50!black,
  fonttitle=\bfseries, title={Phương pháp}, #1%
}
\newtcolorbox{luuy}[1][]{%
  enhanced, breakable,
  colback=red!3!white, colframe=red!50!black,
  fonttitle=\bfseries, title={Lưu ý}, #1%
}
\newtcolorbox{congthuc}[1][]{%
  enhanced, breakable,
  colback=cyan!3!white, colframe=cyan!50!black,
  fonttitle=\bfseries, title={Công thức}, #1%
}
\newtcolorbox{vidu}[1][]{%
  enhanced, breakable,
  colback=violet!3!white, colframe=violet!50!black,
  fonttitle=\bfseries, title={Ví dụ}, #1%
}
\newtcolorbox{phanvidu}[1][]{%
  enhanced, breakable,
  colback=red!2!white, colframe=red!40!black,
  fonttitle=\bfseries, title={Phản ví dụ}, #1%
}
\newtcolorbox{tongket}[1][]{%
  enhanced, breakable,
  colback=blue!4!white, colframe=blue!70!black,
  fonttitle=\bfseries, title={Tổng kết chương}, #1%
}


\usepackage{tabularx}
\usepackage{array}
\usepackage{listings}
\usepackage{float}

\lstset{
  basicstyle=\ttfamily\small,
  breaklines=true,
  frame=single,
  backgroundcolor=\color{gray!8},
  columns=fullflexible,
  keepspaces=true,
}

\graphicspath{{images/intro-to-ml/}}

\newcommand{\mlimg}[2][0.88]{\begin{center}\includegraphics[width=#1\linewidth]{#2}\end{center}}

\begin{document}

\thispagestyle{empty}
\begin{tikzpicture}[remember picture, overlay]
  \fill[blue!80!black] (current page.south west) rectangle (current page.north east);
  \fill[blue!60!cyan, opacity=0.8]
    (current page.south west) -- (current page.north west) --
    ([xshift=-3cm]current page.north east) -- (current page.south east) -- cycle;
  \foreach \r/\o in {6/0.05, 4.5/0.08, 3/0.12} {
    \fill[white, opacity=\o] ([xshift=2cm, yshift=-2cm]current page.north east) circle (\r cm);
  }
  \draw[white, line width=2pt] ([yshift=-5cm]current page.north west) ++(2cm,0) -- ++(17cm,0);
  \draw[white, line width=2pt] ([yshift=-16cm]current page.north west) ++(2cm,0) -- ++(17cm,0);
  \node[anchor=west, white, font=\fontsize{28}{34}\bfseries\selectfont]
    at ([xshift=3cm, yshift=-7.5cm]current page.north west) {\coverbookline};
  \node[anchor=west, white, font=\fontsize{20}{26}\selectfont]
    at ([xshift=3cm, yshift=-9cm]current page.north west) {\covermaintitle};
  \node[anchor=west, white, font=\fontsize{16}{22}\selectfont]
    at ([xshift=3cm, yshift=-10.2cm]current page.north west) {\coversubtitle};
  \node[anchor=west, white, font=\Large\bfseries]
    at ([xshift=3cm, yshift=-13cm]current page.north west) {Biên soạn:};
  \node[anchor=west, yellow!80!white, font=\fontsize{24}{30}\bfseries\selectfont]
    at ([xshift=3cm, yshift=-14.5cm]current page.north west) {ĐẶNG TẤN QUÍ};
  \node[anchor=west, white!90, font=\large]
    at ([xshift=3cm, yshift=-18cm]current page.north west) {SĐT: 0377\,122\,210};
  \node[anchor=south, white!80, font=\small]
    at ([yshift=1.5cm]current page.south) {Tài liệu có bản quyền --- Cấm sao chép khi chưa được sự đồng ý của tác giả};
  \node[anchor=south east, white, font=\Large\bfseries]
    at ([xshift=-2cm, yshift=3cm]current page.south east) {\the\year};
\end{tikzpicture}

\newpage
\tableofcontents
\newpage

\section*{Lời nói đầu}
\addcontentsline{toc}{section}{Lời nói đầu}

\begin{ynghia}[title={Cách đọc tài liệu}]
Biên soạn và dịch từ khóa \emph{Introduction to Machine Learning} của Google Developers\footnote{\url{https://developers.google.com/machine-learning/intro-to-ml}}.
\end{ynghia}

\newpage

%==============================================================================
\section{Học máy là gì?}
%==============================================================================

\begin{dongco}[title={Tại sao cần học máy?}]
Dịch máy, xe tự lái, gợi ý nhạc, tự hoàn thành câu, tóm tắt báo, sinh ảnh chưa từng có: đều cần \textbf{học máy} (ML). ML là cách mới để \emph{giải bài toán}, \emph{trả lời câu hỏi phức tạp} và \emph{tạo nội dung} từ dữ liệu --- thay vì lập trình từng quy tắc bằng tay cho mọi trường hợp.
\end{dongco}

\begin{ungdung}
\begin{itemize}
  \item Dự báo thời tiết, thời gian di chuyển, gợi ý bài hát.
  \item Tự động hoàn thành, tóm tắt văn bản, sinh ảnh/âm thanh/video.
\end{itemize}
\end{ungdung}

\begin{dinhnghia}[title={Học máy}]
ML là quá trình
\href{https://developers.google.com/machine-learning/glossary\#training}{\textbf{huấn luyện}} (training) một phần mềm ---
\href{https://developers.google.com/machine-learning/glossary\#model}{\textbf{mô hình}} (model) ---
để đưa ra
\href{https://developers.google.com/machine-learning/glossary\#prediction}{\textbf{dự đoán}} (prediction) hữu ích hoặc \textbf{sinh nội dung} từ dữ liệu.
\end{dinhnghia}

\begin{ynghia}
\textbf{Mô hình} giống ``công thức ẩn'' học từ dữ liệu: không ai viết sẵn từng quy tắc ``nếu độ ẩm $> 80\%$ thì mưa'', mà máy \emph{tự} tìm quan hệ giữa độ ẩm, áp suất,\ldots\ và lượng mưa sau khi xem hàng nghìn mẫu có nhãn.
\end{ynghia}

\begin{phuongphap}[title={Hai cách giải bài toán dự báo mưa}]
\textbf{Cách truyền thống:} Mô phỏng khí quyển + bề mặt Trái Đất, giải phương trình động lực học chất lỏng --- chính xác về lý thuyật nhưng cực khó triển khai.

\textbf{Cách ML:} Cho mô hình nhiều dữ liệu thời tiết; mô hình học quan hệ toán giữa mẫu thời tiết và lượng mưa; đưa thời tiết hiện tại $\Rightarrow$ dự đoán mưa.

\medskip
Video: \url{https://www.youtube.com/watch?v=RdC7KZ6OFPA}
\end{phuongphap}

\begin{phanvidu}[title={``Mô hình''}]
\begin{itemize}
  \item \textbf{Không} phải phần cứng máy tính.
  \item \textbf{Không} phải mô hình thu nhỏ vật lý (mô hình tàu thủy, mô hình nhà).
  \item \textbf{Phải là:} quan hệ toán học (tập số/tham số) suy ra từ dữ liệu để dự đoán.
\end{itemize}
\end{phanvidu}

\begin{luuy}
Trong ML, ``mô hình'' là gì?

\medskip
Quan hệ toán học học từ dữ liệu, dùng để dự đoán hoặc sinh nội dung.
\end{luuy}

%==============================================================================
\section{Các loại hệ thống học máy}
%==============================================================================

\begin{dongco}[title={Cần phân loại hệ thống ML}]
Cùng là ``học từ dữ liệu'', nhưng bài toán khác nhau: có sẵn đáp án (dự đoán giá nhà), không có đáp án (gom nhóm khách hàng), học qua thử--sai (robot), hay \emph{tạo} nội dung mới (ảnh từ mô tả). Chọn sai loại $\Rightarrow$ sai cách chuẩn bị dữ liệu và đánh giá.
\end{dongco}

\begin{tinhchat}[title={Bốn nhóm hệ thống ML}]
\begin{enumerate}
  \item \textbf{Học có giám sát} --- dữ liệu có nhãn/đáp án đúng.
  \item \textbf{Học không giám sát} --- tìm cấu trúc, cụm, mẫu; không có nhãn cho target.
  \item \textbf{Học tăng cường} --- học qua phần thưởng/hình phạt trong môi trường.
  \item \textbf{AI tạo sinh} --- tạo văn bản, ảnh, âm thanh, video từ đầu vào người dùng.
\end{enumerate}
Một hệ thống có thể kết hợp nhiều nhóm (ví dụ tạo sinh: huấn luyện không giám sát rồi tinh chỉnh có giám sát).
Các loại bổ sung :
\begin{itemize}
  \item Học tự giám sát (Self-Supervised Learning - SLM): Phương pháp này thường được coi là một nhánh nhỏ của học không giám sát, nhưng nó đã phát triển thành một lĩnh vực riêng nhờ thành công trong việc huấn luyện các mô hình quy mô lớn. Nó tự tạo ra nhãn từ dữ liệu mà không cần bất kỳ thao tác gắn nhãn thủ công nào.
  \item Học bán giám sát (Semi-Supervised Learning - SSL): Phương pháp này kết hợp một lượng nhỏ dữ liệu được gán nhãn với một lượng lớn dữ liệu chưa được gán nhãn. Nó hữu ích khi việc gán nhãn dữ liệu tốn kém hoặc mất nhiều thời gian.
\end{itemize}
\end{tinhchat}

\subsection{Học có giám sát}

\begin{dinhnghia}[title={Học có giám sát (supervised learning)}]
\href{https://developers.google.com/machine-learning/glossary\#supervised-machine-learning}{Học có giám sát}: huấn luyện trên dữ liệu có \textbf{đáp án đúng}, tìm mối liên hệ trong dữ liệu dẫn tới đáp án đó. ``Giám sát'' = con người cung cấp kết quả đúng đã biết.
\end{dinhnghia}

\begin{ynghia}
Giống học sinh ôn đề cũ \textbf{có đáp án}: sau đủ đề, làm đề mới tốt hơn. Khác không giám sát: đề \textbf{không} có đáp án, chỉ nhóm câu giống nhau.
\end{ynghia}

\begin{tinhchat}[title={Hai nhánh chính của có giám sát}]
\begin{itemize}
  \item \textbf{Hồi quy} --- đầu ra là \textbf{số} (giá nhà, lượng mưa, thời gian đi).
  \item \textbf{Phân loại} --- đầu ra là \textbf{lớp} / xác suất thuộc lớp (spam, có mèo, loại thời tiết).
\end{itemize}
\end{tinhchat}

\subsubsection{Hồi quy}

\begin{dinhnghia}[title={Mô hình hồi quy}]
\href{https://developers.google.com/machine-learning/glossary\#regression-model}{Mô hình hồi quy} dự đoán một \textbf{giá trị số liên tục} (hoặc số rời có ý nghĩa thứ tự), ví dụ lượng mưa tính bằng mm.
\end{dinhnghia}

\begin{ungdung}[title={Ví dụ hồi quy}]
\begin{tabularx}{\textwidth}{@{}>{\raggedright\arraybackslash}p{2.8cm} X >{\raggedright\arraybackslash}p{3.2cm}@{}}
\toprule
\textbf{Tình huống} & \textbf{Đầu vào} & \textbf{Dự đoán số} \\
\midrule
Giá nhà &
Diện tích, mã ZIP, phòng ngủ/tắm, lãi suất, thuế, số nhà bán khu vực\ldots &
Giá căn \\
Thời gian đi &
Lịch sử giao thông, khoảng cách, thời tiết &
Phút/giây đến nơi \\
\bottomrule
\end{tabularx}
\end{ungdung}

\subsubsection{Phân loại}

\begin{dinhnghia}[title={Mô hình phân loại}]
\href{https://developers.google.com/machine-learning/glossary\#classification-model}{Mô hình phân loại} dự đoán \textbf{khả năng} thuộc một \textbf{lớp} (category), không phải một số tùy ý.
\end{dinhnghia}

\begin{tinhchat}[title={Nhị phân và đa lớp}]
\begin{itemize}
  \item \textbf{Nhị phân:} hai lớp --- \texttt{mưa} / \texttt{không mưa}, \texttt{spam} / \texttt{không spam}.
  \item \textbf{Đa lớp:} $\geq 3$ lớp --- \texttt{mưa}, \texttt{mưa đá}, \texttt{tuyết}, \texttt{mưa phùn}.
\end{itemize}
\end{tinhchat}

\begin{ungdung}
Lọc email spam; nhận diện mèo trong ảnh; phân loại loại thời tiết. Video: \url{https://www.youtube.com/watch?v=NNue6rkDrws}
\end{ungdung}

\begin{phanvidu}[title={Nhầm hồi quy với phân loại}]
Dự đoán \textbf{kWh/năm} là \textbf{hồi quy} (số), không phải phân loại. Phân loại dùng khi cần \textbf{lớp} (spam/không spam), không phải khi nhãn là số đo liên tục trừ khi ta \emph{chủ động} chia khoảng thành lớp (cao/trung bình/thấp).
\end{phanvidu}

\begin{luuy}
Dự đoán tiêu thụ điện tòa nhà (kWh) $\Rightarrow$ \textbf{hồi quy}.
\end{luuy}

\subsection{Học không giám sát}

\begin{dinhnghia}[title={Học không giám sát}]
\href{https://developers.google.com/machine-learning/glossary\#unsupervised-machine-learning}{Học không giám sát} tìm \textbf{mẫu có ý nghĩa} trong dữ liệu \textbf{không có nhãn} cho target. Kỹ thuật phổ biến:
\href{https://developers.google.com/machine-learning/glossary\#clustering}{\textbf{phân cụm}} (clustering) --- gom điểm tương tự thành \textbf{cụm}.
\end{dinhnghia}

\begin{phuongphap}[title={Quy trình phân cụm (trực giác)}]
\begin{enumerate}
  \item Đưa dữ liệu (không cần nhãn target).
  \item Thuật toán đo ``gần nhau'' (khoảng cách, độ tương tự).
  \item Gom thành cụm; \textbf{người} đặt tên cụm sau (mùa, loại khách hàng\ldots).
\end{enumerate}
\end{phuongphap}

\mlimg{clustering-02.png}
\begin{center}\small\textbf{Hình 1.} Phân cụm điểm tương tự.\end{center}

\mlimg{clustering-04.png}
\begin{center}\small\textbf{Hình 2.} Cụm có ranh giới tự nhiên.\end{center}

\begin{ynghia}[title={Phân cụm khác phân loại}]
\textbf{Phân loại:} bạn định nghĩa trước lớp (\texttt{mèo}, \texttt{chó}). \textbf{Phân cụm:} thuật toán tìm nhóm; bạn \emph{đoán} nhãn sau (ví dụ cụm nhiệt độ thấp = ``mùa đông'').
\end{ynghia}

\mlimg{clustering-01.png}
\begin{center}\small\textbf{Hình 3.} Cụm mẫu thời tiết.\end{center}

\mlimg{clustering-03.png}
\begin{center}\small\textbf{Hình 4.} Cụm được gán nhãn tuyết, mưa phùn, mưa, không mưa.\end{center}

\begin{phanvidu}
\textbf{Sai:} ``Có giám sát thường dùng phân cụm.'' --- Phân cụm là \textbf{không giám sát}.

\textbf{Sai:} ``Không giám sát tự biết nghĩa cụm.'' --- Máy chỉ gom nhóm; \textbf{người} đặt tên dựa trên hiểu biết miền.
\end{phanvidu}

\begin{luuy}
\begin{itemize}
  \item Có giám sát: có đáp án đúng trong dữ liệu.
  \item Không giám sát: không có đáp án; tìm nhóm/mẫu.
\end{itemize}
\end{luuy}

\subsection{Học tăng cường}

\begin{dinhnghia}[title={Học tăng cường (RL)}]
\href{https://developers.google.com/machine-learning/glossary\#reinforcement-learning-rl}{Học tăng cường}: tác nhân thực hiện \textbf{hành động} trong \textbf{môi trường}, nhận
\href{https://developers.google.com/machine-learning/glossary\#reward}{\textbf{phần thưởng}} hoặc \textbf{hình phạt}, học
\href{https://developers.google.com/machine-learning/glossary\#policy}{\textbf{chính sách}} (policy) tối đa hóa phần thưởng tích lũy.
\end{dinhnghia}

\begin{congthuc}[title={target RL (ý tưởng)}]
Tìm chính sách $\pi$ sao cho tổng phần thưởng kỳ vọng lớn:
\[
  \pi^\ast \approx \arg\max_{\pi}\; \mathbb{E}\Big[\sum_t \gamma^t\, r_t\Big],
\]
với $r_t$ phần thưởng bước $t$, $\gamma \in (0,1]$ hệ số chiết khấu (chi tiết ở khóa RL chuyên sâu).
\end{congthuc}

\begin{ungdung}
Huấn luyện robot di chuyển; \href{https://deepmind.com/research/case-studies/alphago-the-story-so-far}{AlphaGo} chơi cờ vây; game, điều khiển tài nguyên.
\end{ungdung}

\subsection{AI tạo sinh}

\begin{dinhnghia}[title={AI tạo sinh (generative AI)}]
\href{https://developers.google.com/machine-learning/glossary\#generative-ai}{AI tạo sinh}: lớp mô hình \textbf{tạo nội dung} từ đầu vào người dùng --- văn bản, ảnh, nhạc, video; hoặc kết hợp (ảnh + chữ $\to$ video).
\end{dinhnghia}

\begin{tinhchat}[title={Ký hiệu kiểu đầu vào $\to$ đầu ra}]
text-to-text, text-to-image, text-to-video, text-to-code, text-to-speech, image-and-text-to-image,\ldots
\end{tinhchat}

\begin{table}[H]
\centering
\caption{Bảng ví dụ mô hình tạo sinh (rút gọn)}
\small
\begin{tabularx}{\textwidth}{@{}>{\raggedright\arraybackslash}p{2.4cm} >{\raggedright\arraybackslash}p{4.2cm} X@{}}
\toprule
\textbf{Mô hình} & \textbf{Đầu vào} & \textbf{Đầu ra} \\
\midrule
Text-to-text & Ai sáng lập đua Le Mans? & Giải thích lịch sử (Bard). \\
Text-to-image & Bạch tuộc ngoài hành tinh đọc báo qua cổng. & Ảnh (Imagen). \\
Text-to-code & Vòng lặp Python in số nguyên tố. & Mã (Bard). \\
Image-to-text & Ảnh hồng hạc & Mô tả văn bản. \\
\bottomrule
\end{tabularx}
\end{table}

\begin{vidu}[title={Text-to-video (Phenaki)}]
\textbf{Đầu vào:} ``Gấu bông bơi ở biển San Francisco\ldots''
\mlimg[0.5]{teddy_bear.png}
Nguồn: \href{https://phenaki.video/}{Phenaki}.
\end{vidu}

\begin{vidu}[title={Text-to-code (Bard)}]
\begin{lstlisting}[language=Python]
for number in numbers:
    is_prime = True
    for i in range(2, number):
        if number % i == 0:
            is_prime = False
            break
    if is_prime:
        print(number)
\end{lstlisting}
\end{vidu}

\begin{vidu}[title={Image-to-text}]
\mlimg[0.35]{flamingo.png}
``Đây là hồng hạc, sống ở vùng Caribbean.''
\end{vidu}

\begin{ynghia}[title={Mô hình tạo sinh học như thế nào?}]
Học \textbf{mẫu} trong dữ liệu để tạo dữ liệu \textbf{mới tương tự} --- như diễn viên bắt chước phong cách, họa sĩ học trường phái, ban cover học nhóm nhạc. Thường: huấn luyện \textbf{không giám sát} ban đầu, sau đó tinh chỉnh \textbf{có giám sát} hoặc \textbf{tăng cường}.
\end{ynghia}

\begin{ungdung}[title={AI tạo sinh trong thương mại điện tử}]
Tự xóa nền ảnh sản phẩm; nâng chất lượng ảnh độ phân giải thấp --- use case đang phát triển nhanh.
\end{ungdung}

\begin{dinhly}[title={Nguyên lý chung của ML (tóm tắt phần 2)}]
\begin{enumerate}
  \item \textbf{Có giám sát} cần nhãn; phù hợp dự đoán khi biết ``đáp án'' cần học.
  \item \textbf{Không giám sát} khám phá cấu trúc; nhãn cụm do người gán sau.
  \item \textbf{Tăng cường} học hành vi qua phần thưởng.
  \item \textbf{tạo sinh} tạo nội dung mới từ mẫu đã học.
\end{enumerate}
\end{dinhly}

\begin{tongket}[title={Thuật ngữ --- Các loại hệ thống ML}]
\href{https://developers.google.com/machine-learning/glossary\#classification-model}{classification model} \textbullet\
\href{https://developers.google.com/machine-learning/glossary\#clustering}{clustering} \textbullet\
\href{https://developers.google.com/machine-learning/glossary\#model}{model} \textbullet\
\href{https://developers.google.com/machine-learning/glossary\#policy}{policy} \textbullet\
\href{https://developers.google.com/machine-learning/glossary\#prediction}{prediction} \textbullet\
\href{https://developers.google.com/machine-learning/glossary\#regression-model}{regression model} \textbullet\
\href{https://developers.google.com/machine-learning/glossary\#reinforcement-learning-rl}{reinforcement learning} \textbullet\
\href{https://developers.google.com/machine-learning/glossary\#reward}{reward} \textbullet\
\href{https://developers.google.com/machine-learning/glossary\#supervised-machine-learning}{supervised learning} \textbullet\
\href{https://developers.google.com/machine-learning/glossary\#training}{training} \textbullet\
\href{https://developers.google.com/machine-learning/glossary\#unsupervised-machine-learning}{unsupervised learning}.
\end{tongket}

\newpage

%==============================================================================
\section{Học có giám sát --- Khái niệm nền}
%==============================================================================

\begin{dongco}[title={Vì sao cần năm khái niệm?}]
Lọc spam hay dự báo mưa đều cùng ``khung'': có \textbf{dữ liệu có nhãn}, xây \textbf{mô hình}, \textbf{huấn luyện}, \textbf{đánh giá}, rồi \textbf{suy luận} trên dữ liệu mới. Thiếu một bước (ví dụ đánh giá kém mà triển khai) $\Rightarrow$ sản phẩm sai trong thực tế.
\end{dongco}

\begin{phuongphap}[title={Vòng đời mô hình có giám sát}]
\begin{center}
\textbf{Dữ liệu} $\longrightarrow$ \textbf{Mô hình} $\longrightarrow$ \textbf{Huấn luyện} $\longrightarrow$ \textbf{Đánh giá} $\longrightarrow$ \textbf{Suy luận}
\end{center}
Có thể lặp Huấn luyện $\leftrightarrow$ Đánh giá nhiều lần trước khi suy luận trong sản phẩm.
\end{phuongphap}

\subsection{Dữ liệu}

\begin{dinhnghia}[title={Tập dữ liệu, mẫu, đặc trưng, nhãn}]
\textbf{Tập dữ liệu} (dataset): tập dữ liệu liên quan (ảnh mèo, giá nhà, thời tiết).

\textbf{Mẫu} (\href{https://developers.google.com/machine-learning/glossary\#example}{example}): một đơn vị --- như một đường thẳng bảng.

\textbf{Đặc trưng} (\href{https://developers.google.com/machine-learning/glossary\#feature}{feature}): giá trị đầu vào mô hình dùng để dự đoán.

\textbf{Nhãn} (\href{https://developers.google.com/machine-learning/glossary\#label}{label}): ``đáp án'' cần dự đoán.

\textbf{Mẫu có nhãn} (\href{https://developers.google.com/machine-learning/glossary\#labeled-example}{labeled example}): có cả đặc trưng và nhãn.
\end{dinhnghia}

\begin{vidu}[title={Mô hình dự báo mưa}]
Đặc trưng: vĩ độ, kinh độ, nhiệt độ, độ ẩm, mây, gió, áp suất. Nhãn: \emph{lượng mưa}.
\end{vidu}

\begin{tinhchat}[title={Mẫu có nhãn và không nhãn}]
\begin{itemize}
  \item \textbf{Có nhãn:} dùng khi \textbf{huấn luyện} và \textbf{đánh giá} (so sánh dự đoán với nhãn thật).
  \item \textbf{Không nhãn:} dùng khi \textbf{suy luận} trên dữ liệu mới --- mô hình tự dự đoán nhãn.
\end{itemize}
\end{tinhchat}

\paragraph{Hai mẫu có nhãn.}
\mlimg{labeled_example.png}

\paragraph{Hai mẫu không nhãn.}
\mlimg{unlabeled_example.png}

\begin{tinhchat}[title={Kích thước và đa dạng tập dữ liệu}]
\textbf{Kích thước:} số mẫu. \textbf{Đa dạng:} phạm vi tình huống (mùa, khu vực, loại nhà\ldots). Tập lý tưởng: \textbf{lớn và đa dạng}.
\end{tinhchat}

\begin{phanvidu}[title={Tập lớn nhưng không đa dạng}]
100 năm dữ liệu \textbf{chỉ tháng 7} $\Rightarrow$ dự báo mưa tháng 1 kém dù ``nhiều dữ liệu''.
\end{phanvidu}

\begin{phanvidu}[title={Tập đa dạng nhưng quá nhỏ}]
Vài năm, đủ 12 tháng, nhưng ít năm $\Rightarrow$ thiếu biến thiên dài hạn; dự đoán thiếu tin cậy.
\end{phanvidu}

\begin{luuy}[title={Tập dữ liệu lý tưởng}]
\textbf{Đáp án:} kích thước \textbf{lớn} và đa dạng \textbf{cao}. Tránh: lớn/thấp đa dạng; nhỏ/cao đa dạng; nhỏ/thấp đa dạng.
\end{luuy}

\begin{luuy}[title={Nhiều đặc trưng không luôn tốt hơn}]
Hàng trăm đặc trưng (ảnh vệ tinh, mây\ldots) có thể giúp, nhưng đặc trưng \textbf{không nhân quả} với nhãn (ví dụ màu xe với giá) có thể \textbf{không} cải thiện dự đoán.
\end{luuy}

\subsection{Mô hình}

\begin{dinhnghia}[title={Mô hình trong học có giám sát}]
Tập hợp số (tham số) định nghĩa quan hệ toán từ \emph{mẫu đặc trưng} tới \emph{giá trị nhãn}. Quan hệ được \textbf{học} qua huấn luyện, không gán tay đầy đủ.
\end{dinhnghia}

\begin{ynghia}
Có thể hình dung mô hình như hàm $f$: nhãn $\approx f(\text{đặc trưng}_1, \ldots, \text{đặc trưng}_n)$; các hệ số trong $f$ được chỉnh khi huấn luyện.
\end{ynghia}

\subsection{Huấn luyện}

\begin{dinhnghia}[title={Huấn luyện (training)}]
Cho mô hình tập \textbf{mẫu có nhãn}; mô hình dự đoán, so với nhãn thật, cập nhật nghiệm để dự đoán tốt hơn trung bình trên tập.
\end{dinhnghia}

\begin{congthuc}[title={Mất mát (loss) --- ý tưởng}]
Chênh lệch giữa dự đoán $\hat{y}$ và nhãn thật $y$:
\[
  \text{loss} = L(\hat{y}, y).
\]
Huấn luyện $\Leftrightarrow$ giảm loss trung bình trên các mẫu (cập nhật tham số mô hình).
\end{congthuc}

\begin{vidu}[title={Một bước cập nhật}]
Dự đoán \texttt{1{,}15} inch mưa, thực tế \texttt{0{,}75} inch $\Rightarrow$ chỉnh mô hình để lần sau gần \texttt{0{,}75} hơn.
\end{vidu}

\begin{phuongphap}[title={Ba bước huấn luyện (theo Google ML)}]
\begin{enumerate}
  \item Một mẫu có nhãn $\to$ dự đoán.
  \mlimg[0.72]{training-a-model-01.png}
  \begin{center}\small\textbf{Hình 5.} Dự đoán từ mẫu có nhãn.\end{center}

  \item So sánh với nhãn thật $\to$ cập nhật.
  \mlimg[0.72]{training-a-model-02.png}
  \begin{center}\small\textbf{Hình 6.} Cập nhật dự đoán.\end{center}

  \item Lặp cho mọi mẫu (có thể nhiều \textbf{epoch}).
  \mlimg[0.72]{training-a-model-03.png}
  \begin{center}\small\textbf{Hình 7.} Cập nhật trên cả tập.\end{center}
\end{enumerate}
\end{phuongphap}

\begin{luuy}[title={Chọn đặc trưng khi huấn luyện}]
Thử thêm/bỏ đặc trưng (ví dụ \texttt{time\_of\_day} trong dữ liệu thời tiết) và xem đánh giá có cải thiện --- không phải đặc trưng nào cũng có \textbf{sức dự đoán}.
\end{luuy}

\subsection{Đánh giá}

\begin{phuongphap}[title={Đánh giá mô hình đã huấn luyện}]
\begin{enumerate}
  \item Dùng tập có nhãn (thường \textbf{tách} khỏi tập huấn luyện).
  \item Chỉ đưa \textbf{đặc trưng} cho mô hình.
  \item So sánh dự đoán với nhãn thật.
  \item Nếu chưa đạt $\Rightarrow$ huấn luyện lại / đổi đặc trưng / thu thập thêm dữ liệu.
\end{enumerate}
\end{phuongphap}

\mlimg[0.8]{evaluating-a-model.png}
\begin{center}\small\textbf{Hình 8.} So dự đoán với giá trị thực.\end{center}

\begin{luuy}[title={Vì sao phải huấn luyện?}]
Mô hình phải \textbf{học} quan hệ đặc trưng--nhãn từ dữ liệu; không ``biết sẵn'' khi mới khởi tạo.
\end{luuy}

\subsection{Suy luận}

\begin{dinhnghia}[title={Suy luận (inference)}]
\href{https://developers.google.com/machine-learning/glossary\#inference}{Suy luận}: dùng mô hình đã huấn luyện và đánh giá đủ tốt để dự đoán trên mẫu \textbf{không có nhãn} (chỉ có đặc trưng) --- ví dụ thời tiết hiện tại $\to$ lượng mưa dự kiến.
\end{dinhnghia}

\begin{dinhly}[title={Chuỗi logic có giám sát}]
Huấn luyện chỉ trên mẫu \textbf{có nhãn} $\Rightarrow$ đánh giá trên mẫu có nhãn (ẩn nhãn với mô hình) $\Rightarrow$ suy luận trên mẫu \textbf{không nhãn} trong vận hành thực tế.
\end{dinhly}

\begin{tongket}[title={Thuật ngữ --- Học có giám sát}]
\href{https://developers.google.com/machine-learning/glossary\#example}{example} \textbullet\
\href{https://developers.google.com/machine-learning/glossary\#feature}{feature} \textbullet\
\href{https://developers.google.com/machine-learning/glossary\#inference}{inference} \textbullet\
\href{https://developers.google.com/machine-learning/glossary\#labeled-example}{labeled example} \textbullet\
\href{https://developers.google.com/machine-learning/glossary\#label}{label} \textbullet\
\href{https://developers.google.com/machine-learning/glossary\#loss}{loss} \textbullet\
\href{https://developers.google.com/machine-learning/glossary\#prediction}{prediction} \textbullet\
\href{https://developers.google.com/machine-learning/glossary\#training}{training}.
\end{tongket}

\newpage

%==============================================================================
\section{Kiểm tra hiểu biết tổng hợp}
%==============================================================================

\begin{dongco}[title={Luyện chọn mô hình và đặc trưng}]
Biết định nghĩa chưa đủ: cần nhìn bảng dữ liệu và hỏi --- nhãn là số hay lớp? Có nhãn trong tập không? Đặc trưng nào thực sự liên quan giá/nhãn?
\end{dongco}

\subsection{Sức mạnh dự đoán}

\begin{ynghia}[title={Predictive power}]
Không phải mọi cột trong bảng đều giúp dự đoán nhãn. Chỉ vài đặc trưng mang phần lớn thông tin (ví dụ hãng xe, năm, số dặm với giá).
\end{ynghia}

\mlimg{predictive_power.png}

\begin{luuy}[title={Dự đoán giá xe}]
\textbf{Tốt nhất:} \texttt{make\_model}, \texttt{year}, \texttt{miles}.

\textbf{Yếu:} \texttt{color}, \texttt{height}; \texttt{gearbox}; \texttt{tire\_size}, \texttt{wheel\_base}.
\end{luuy}

\subsection{Có giám sát và không giám sát}

\begin{vidu}[title={Người dùng website mua sắm}]
\mlimg[0.95]{supervised-unsupervised.png}

Muốn \textbf{phân loại kiểu người dùng} (chưa có nhãn) $\Rightarrow$ \textbf{không giám sát} (phân cụm), sau đó đặt tên: ``săn giảm giá'', ``deal hunter'', ``surfer'', ``loyal'', ``wanderer''.
\end{vidu}

\begin{phanvidu}[title={Sai: dùng có giám sát khi không có nhãn}]
Tập không có cột ``loại người dùng'' $\Rightarrow$ không thể huấn luyện có giám sát để ``dự đoán lớp người dùng'' trừ khi tự tạo nhãn (gán nhãn tay).
\end{phanvidu}

\begin{vidu}[title={Tiêu thụ điện hộ gia đình}]
\mlimg[0.95]{supervised.png}

Dự đoán \textbf{kWh/năm} $\Rightarrow$ \textbf{có giám sát}; nhãn = \texttt{kilowatt hours used per year}; đặc trưng = diện tích, vị trí, năm xây.
\end{vidu}

\begin{vidu}[title={Giá vé máy bay}]
\mlimg[0.95]{flight-table.png}

\textbf{Hồi quy} (giá vé là số). Có thể \textbf{phân loại} ``cao/trung bình/thấp'' nếu chuyển số thành lớp theo ngưỡng (sau khi tính trung bình tuyến).
\end{vidu}

\begin{tinhchat}[title={Hồi quy vs phân loại cho giá vé}]
\begin{itemize}
  \item Hồi quy: đầu ra số (USD).
  \item Phân loại đa lớp: có thể 3 lớp \texttt{cao}, \texttt{trung bình}, \texttt{thấp} --- cần tiền xử lý nhãn.
\end{itemize}
\end{tinhchat}

\subsection{Huấn luyện và đánh giá}

\begin{phuongphap}[title={Khi dự đoán lệch nhiều --- hai hướng hiệu quả}]
\begin{enumerate}
  \item Huấn luyện lại với \textbf{ít đặc trưng hơn} nhưng có sức dự đoán mạnh.
  \item Huấn luyện lại trên tập \textbf{lớn hơn, đa dạng hơn}.
\end{enumerate}
\end{phuongphap}

\begin{phanvidu}
\textbf{Sai:} ``Không sửa được mô hình dự đoán tệ.'' --- Có thể: thêm dữ liệu, đổi đặc trưng, nhiều vòng huấn luyện.

\textbf{Sai:} ``Đổi sang không giám sát sẽ cải thiện dự đoán có nhãn.'' --- Không giải quyết bài toán dự đoán nhãn đã biết trong huấn luyện.
\end{phanvidu}

\begin{luuy}[title={Cải thiện mô hình}]
Chọn \textbf{hai} đáp án tốt: (1) huấn luyện lại với đặc trưng mạnh; (2) tập lớn/đa dạng hơn. Không chọn: ``không sửa được''; ``đổi sang không giám sát''.
\end{luuy}

\begin{tongket}[title={Bước tiếp theo trong hành trình ML}]
\begin{itemize}
  \item \href{https://developers.google.com/machine-learning/crash-course}{Machine Learning Crash Course} --- thực hành sâu.
  \item \href{https://developers.google.com/machine-learning/problem-framing}{Problem Framing} --- đặt khung bài toán, tránh lỗi.
  \item \href{https://pair.withgoogle.com/guidebook/}{People + AI Guidebook} --- thiết kế AI lấy con người làm trung tâm.
\end{itemize}
\end{tongket}

\end{document}
